% Compile with: uplatex file.tex && dvipdfmx file.dvi
%           or: latexmk -uplatex file.tex
%\documentclass[uplatex,dvipdfmx]{jsarticle}
\documentclass[dvipdfmx,a4paper,11pt]{jsarticle}
\usepackage{amsmath, amssymb}
\usepackage{tikz}

\title{scenario\_maxwell\_consistency}
\date{\today}

\begin{document}
\maketitle

電磁場を扱う。電場 $\mathbf E$ と磁場 $\mathbf B$ をベクトル場として準備する（成分は抽象のままで登録し、後で具体的な式を割り当てる）。

E をベクトル場として定義する

\[
  E
\]

B をベクトル場として定義する

\[
  B
\]

$+x$ 方向に進む直線偏波の平面波として、電場の各成分（振幅1、波数 $k$、角振動数 $\omega$）を割り当てる。

$E_x = 0$

\begin{equation}
  E_x = 0
  \label{eq:E1}
\end{equation}

$E_y = cos(k x - w t)$

\begin{equation}
  E_y = \cos(k x - w t)
  \label{eq:E2}
\end{equation}

$E_z = 0$

\begin{equation}
  E_z = 0
  \label{eq:E3}
\end{equation}

対応する磁場の各成分を割り当てる。

$B_x = 0$

\begin{equation}
  B_x = 0
  \label{eq:E4}
\end{equation}

$B_y = 0$

\begin{equation}
  B_y = 0
  \label{eq:E5}
\end{equation}

$B_z = cos(k x - w t)/c$

\begin{equation}
  B_z = \cos(k x - w t)/c
  \label{eq:E6}
\end{equation}

電場の発散 $\nabla\cdot\mathbf E$ を計算する。真空中（電荷なし）では恒等的にゼロになるはずである（Gaussの法則）。

$\nabla \cdot E$

\begin{equation}
  0
  \label{eq:E7}
\end{equation}

磁場の発散 $\nabla\cdot\mathbf B$ を計算する。磁荷が存在しないため恒等的にゼロになるはずである。

$\nabla \cdot B$

\begin{equation}
  0
  \label{eq:E8}
\end{equation}

電場の回転 $\nabla\times\mathbf E$ を計算する（Faradayの法則の左辺）。

$\nabla \times E$

\begin{equation}
  \left[ 0 , 0 , -\left(k \sin \left(k x-t w\right)\right) \right]
  \label{eq:E9}
\end{equation}

磁場の回転 $\nabla\times\mathbf B$ を計算する（Ampèreの法則の左辺）。

$\nabla \times B$

\begin{equation}
  \left[ 0 , \frac{k \sin \left(k x-t w\right)}{c} , 0 \right]
  \label{eq:E10}
\end{equation}

$\partial/\partial t$ はベクトル演算子として直接扱えないため、成分ごとに時間偏微分する。ここではゼロでない成分 $B_z$ を $t$ で偏微分する。

$B_z$ を t で偏微分して

\begin{equation}
  \frac{\partial}{\partial t} B_z(x,y,z) = \frac{w \sin \left(k x-t w\right)}{c}
  \label{eq:E11}
\end{equation}

求めた偏微分の結果を、新しいベクトル場として組み立てる。ここで `dBdt` という名前を付けるが、これは $d\mathbf B/dt$ を表すために選んだ単なる識別子であり（CASの予約語や特別な構文ではない、他の名前でも同様に動く）、成分ごとに割り当て直す。

dBdt をベクトル場として定義する

\[
  dBdt
\]

$dBdt_x = 0$

\begin{equation}
  dBdt_x = 0
  \label{eq:E12}
\end{equation}

$dBdt_y = 0$

\begin{equation}
  dBdt_y = 0
  \label{eq:E13}
\end{equation}

$dBdt_z = w sin(k x - w t)/c$

\begin{equation}
  dBdt_z = w \sin(k x - w t)/c
  \label{eq:E14}
\end{equation}

組み立てた $d\mathbf B/dt$ が後続の計算（発散など）に使えることを、その発散を計算して確認する。

$\nabla \cdot dBdt$

\begin{equation}
  0
  \label{eq:E15}
\end{equation}

Faradayの法則 $\nabla\times\mathbf E=-\partial\mathbf B/\partial t$ の $z$ 成分どうしを等式として登録する（両辺の $z$ 成分を比較する）。

$-k sin(k x - w t) = -w sin(k x - w t)/c$

\begin{equation}
  -k \sin(k x - w t) = -w \sin(k x - w t)/c
  \label{eq:E16}
\end{equation}

この等式を角振動数 $\omega$ について解き、整合条件（分散関係）$\omega=ck$ を導く。

E16 を w について解いて

\begin{equation}
  w = \frac{k x}{t}
  \label{eq:E17}
\end{equation}

\begin{equation}
  w = c k
  \label{eq:E18}
\end{equation}

\end{document}